%-------------------------
% Resume in Latex
% Author : Jake Gutierrez
% Based off of: https://github.com/sb2nov/resume
% License : MIT
%------------------------

\documentclass[letterpaper,11pt]{article}

\usepackage{latexsym}
\usepackage[empty]{fullpage}
\usepackage{titlesec}
\usepackage{marvosym}
\usepackage[usenames,dvipsnames]{color}
\usepackage{verbatim}
\usepackage{enumitem}
\usepackage[hidelinks]{hyperref}
\usepackage{fancyhdr}
\usepackage[english]{babel}
\usepackage{tabularx}
\input{glyphtounicode}


%----------FONT OPTIONS----------
% sans-serif
% \usepackage[sfdefault]{FiraSans}
% \usepackage[sfdefault]{roboto}
% \usepackage[sfdefault]{noto-sans}
% \usepackage[default]{sourcesanspro}

% serif
% \usepackage{CormorantGaramond}
% \usepackage{charter}


\pagestyle{fancy}
\fancyhf{} % clear all header and footer fields
\fancyfoot{}
\renewcommand{\headrulewidth}{0pt}
\renewcommand{\footrulewidth}{0pt}

% Adjust margins
\addtolength{\oddsidemargin}{-0.5in}
\addtolength{\evensidemargin}{-0.5in}
\addtolength{\textwidth}{1in}
\addtolength{\topmargin}{-.5in}
\addtolength{\textheight}{1.0in}

\urlstyle{same}

\raggedbottom
\raggedright
\setlength{\tabcolsep}{0in}

% Sections formatting
\titleformat{\section}{
  \vspace{-4pt}\scshape\raggedright\large
}{}{0em}{}[\color{black}\titlerule \vspace{-5pt}]

% Ensure that generate pdf is machine readable/ATS parsable
\pdfgentounicode=1

%-------------------------
% Custom commands
\newcommand{\resumeItem}[1]{
  \item\small{
    {#1 \vspace{-2pt}}
  }
}

\newcommand{\resumeSubheading}[4]{
  \vspace{-2pt}\item
    \begin{tabular*}{0.97\textwidth}[t]{l@{\extracolsep{\fill}}r}
      \textbf{#1} & #2 \\
      \textit{\small#3} & \textit{\small #4} \\
    \end{tabular*}\vspace{-7pt}
}

\newcommand{\resumeSubSubheading}[2]{
    \item
    \begin{tabular*}{0.97\textwidth}{l@{\extracolsep{\fill}}r}
      \textit{\small#1} & \textit{\small #2} \\
    \end{tabular*}\vspace{-7pt}
}

\newcommand{\resumeProjectHeading}[2]{
    \item
    \begin{tabular*}{0.97\textwidth}{l@{\extracolsep{\fill}}r}
      \small#1 & #2 \\
    \end{tabular*}\vspace{-7pt}
}

\newcommand{\resumeSubItem}[1]{\resumeItem{#1}\vspace{-4pt}}

\renewcommand\labelitemii{$\vcenter{\hbox{\tiny$\bullet$}}$}

\newcommand{\resumeSubHeadingListStart}{\begin{itemize}[leftmargin=0.15in, label={}]}
\newcommand{\resumeSubHeadingListEnd}{\end{itemize}}
\newcommand{\resumeItemListStart}{\begin{itemize}}
\newcommand{\resumeItemListEnd}{\end{itemize}\vspace{-5pt}}

%-------------------------------------------
%%%%%%  RESUME STARTS HERE  %%%%%%%%%%%%%%%%%%%%%%%%%%%%


\begin{document}

%----------HEADING----------
% \begin{tabular*}{\textwidth}{l@{\extracolsep{\fill}}r}
%   \textbf{\href{http://sourabhbajaj.com/}{\Large Sourabh Bajaj}} & Email : \href{mailto:sourabh@sourabhbajaj.com}{sourabh@sourabhbajaj.com}\\
%   \href{http://sourabhbajaj.com/}{http://www.sourabhbajaj.com} & Mobile : +1-123-456-7890 \\
% \end{tabular*}

\begin{center}
    \textbf{\Huge \scshape Carlos Chavarria} \\ \vspace{1pt}
    \small 999-420-1723 $|$ \href{mailto:cjchavarriamendoza@gmail.com}{\underline{cjchavarriamendoza@gmail.com}} $|$ 
    \href{https://linkedin.com/in/cjchavarria}{\underline{linkedin.com/in/cjchavarria}} $|$
    \href{https://github.com/cj-chavarria}{\underline{github.com/cj-chavarria}}
\end{center}


%-----------EDUCATION-----------
\section{Educación}
  \resumeSubHeadingListStart
    \resumeSubheading
      {Universidad Autónoma de Yucatán}{Mérida, Yucatán, MX}
      {Licenciatura en Ingeniería Física}{Agosto. 2021 -- Mayo 2026}
      \resumeItemListStart
      \resumeItem{Asignaturas Relevantes:}
        \begin{itemize}
            \item Métodos Numéricos, Probabilidad y Estadística.
            \item Cursos optativos: Machine Learning, Ciencia de datos, Deep Learning para Visión por Computadora, Modern Cloud Engineering
        \end{itemize}
      \resumeItemListEnd
    
  \resumeSubHeadingListEnd


%-----------EXPERIENCE-----------
\section{Experience}
  \resumeSubHeadingListStart

    \resumeSubheading
      {Ingeniero en Automatización}{Enero 2026 -- Presente}
      {WAYAKNA S.A. de C.V.}{Mérida, Yucatán, MX}
      \resumeItemListStart
        \resumeItem{Diseñar e implementar un proceso ETL para migrar datos a un software de CRM utilizando Snowflake, Apache Airflow y dbt.}
        \resumeItem{Desarrollé un API Wrapper para diversas API oficiales de registros mercantiles gubernamentales de distintos países, con el fin de facilitar el enriquecimiento y la validación de datos para un sistema de procesamiento de datos}
        \resumeItem{Desarrollé una herramienta de procesamiento de facturas para la extracción de texto utilizando bibliotecas de OCR en Python, que se integró en un proceso de validación de información.}
        \resumeItem{Desarrollé un script de web scraping como parte de un proceso para generar información (correos electrónicos, números de teléfono, etc.) sobre clientes potenciales para el equipo de ventas, lo que permitió obtener con éxito más de 5,000 nuevos clientes potenciales de diversas fuentes.}
        
      \resumeItemListEnd

    \resumeSubheading
    {Especialista QA}{Febrero 2025 -- Enero 2026}
    {Scale AI}{Remote}
        \resumeItemListStart
          \resumeItem{Me especialicé en garantizar la calidad y la precisión de los modelos de inteligencia artificial mediante rigurosos procesos de validación y control. Mi trabajo consistía en la revisión de datos, el análisis del desempeño y las pruebas de calidad.}
        \resumeItemListEnd

  \resumeSubHeadingListEnd

%-----------PROGRAMMING SKILLS-----------
\section{Habilidades Técnicas | Idiomas}
 \begin{itemize}[leftmargin=0.15in, label={}]
    \small{\item{
     \textbf{Lenguajes (Programación)}{: Python, SQL (Postgres)} \\
     \textbf{Herramientas de desarrollo}{: Git, Docker, Google Cloud Platform, Jupyter/Marimo Notebooks.} \\
     \textbf{Librerías}{: polars, pandas, NumPy, Matplotlib, seaborn, scikit-learn, OpenCV, Pydantic, Apache Airflow, Prefect.}
     \textbf{Idioma: }{Ingles (C2 - Oxford Placement Test)} 
    }}
 \end{itemize}



%-----------PROJECTS-----------
\section{Proyectos}
    \resumeSubHeadingListStart
      \resumeProjectHeading
          {\textbf{Análisis de las cafeterías de Mérida, Yuc — Proyecto personal} $|$ \emph{Python}}{}
          \resumeItemListStart
            \resumeItem{Diseñé e implementé un proceso de ingestión de datos en Python utilizando la API de Google Places y cuadrículas geoespaciales para recopilar de manera sistemática información de todas las cafeterías de la ciudad.}
            \resumeItem{Llevé a cabo la limpieza de datos y la detección de anomalías para analizar patrones de negocio basados en la ubicación, los precios y las calificaciones.}
            \resumeItem{Desarrollé visualizaciones geoespaciales interactivas con Folium para representar en mapas la densidad de empresas y los indicadores de desempeño.}
          \resumeItemListEnd
      \resumeProjectHeading
          {\textbf{Análisis de la Calidad del Aire — Proyecto de Investigación} $|$ \emph{Python}}{Enero 2025 -- Agosto 2025}
          \resumeItem{Proyecto: “Medición de la contaminación ambiental al interior de viviendas donde se cocina con leña y carbón en Pixyá, Yucatán”}
          \resumeItemListStart
            \resumeItem{Desarrollé scripts en Python para el preprocesamiento automatizado de datos, la detección de valores atípicos y el análisis de series temporales de indicadores de contaminación ambiental con pandas y Matplotlib}
            \resumeItem{Generar informes estadísticos para el equipo de investigación con el fin de facilitar el análisis de los resultados obtenidos.}
          \resumeItemListEnd
    \resumeSubHeadingListEnd



%

\end{document}
